\section*{Capítulo \ref{ch:g12:curly-arrows} --- Mecanismos de reação: as setas curvas}

\begin{solution}{exo:g12:curly-arrows:1}
\ce{OH-}: doador (\omterm{def:g10:lewis-and-shape:lone-pair}{pares não ligantes}, carga negativa); nenhum sítio
aceptor. \ce{H+}: aceptor (nenhum \omterm{def:g9:inside-the-atom:nucleus}{elétron}). \ce{CH2=CH2}: doador (a
\omterm{def:g10:lewis-and-shape:multiple-bond}{ligação dupla}). \ce{CH3-Cl}: doador (\omterm{def:g10:lewis-and-shape:lone-pair}{pares não ligantes} do \ce{Cl}),
aceptor (o carbono, $\delta^+$). \ce{H2O}: doador (\omterm{def:g10:lewis-and-shape:lone-pair}{pares não ligantes} do
\ce{O}); os seus \omterm{def:g7:atoms-and-molecules:atom}{átomos} de hidrogênio, $\delta^+$, são sítios aceptores
fracos.
\end{solution}

\begin{solution}{exo:g12:curly-arrows:2}
(a) \omterm{def:g12:curly-arrows:categories}{substituição}; (b) \omterm{def:g12:curly-arrows:categories}{adição}; (c) \omterm{def:g12:curly-arrows:categories}{eliminação}; (d) \omterm{def:g12:curly-arrows:categories}{adição}.
\end{solution}

\begin{solution}{exo:g12:curly-arrows:3}
Ela parte de um par de \omterm{def:g9:inside-the-atom:nucleus}{elétrons} (um \omterm{def:g10:lewis-and-shape:lone-pair}{par não ligante} ou uma ligação) e
termina no \omterm{def:g7:atoms-and-molecules:atom}{átomo} para onde o par vai. De um \omterm{def:g10:lewis-and-shape:lone-pair}{par não ligante} do oxigênio
até um carbono: o par vira uma nova ligação \ce{O-C}.
\end{solution}

\begin{solution}{exo:g12:curly-arrows:4}
A amônia cede o \omterm{def:g10:lewis-and-shape:lone-pair}{par não ligante} do seu nitrogênio; a seta vai desse par
até o \ce{H+}. O produto \ce{NH4+} tem a carga $+1$.
\end{solution}

\begin{solution}{exo:g12:curly-arrows:5}
Uma espécie formada numa etapa e consumida numa etapa posterior; ela não
está na equação global. Aqui, o carbocátion \ce{CH3-CH2+}.
\end{solution}

\begin{solution}{exo:g12:curly-arrows:6}
Uma seta de um \omterm{def:g10:lewis-and-shape:lone-pair}{par não ligante} do oxigênio do \ce{OH-} até o carbono;
uma seta da ligação \ce{C-Br} até o bromo. Cargas: $-1$ à esquerda
(\ce{OH-}); à direita, o metanol é neutro e o \ce{Br-} tem $-1$: o
total é $-1$ dos dois lados.
\end{solution}

\begin{solution}{exo:g12:curly-arrows:7}
O bromo é mais eletronegativo que o hidrogênio: ele fica com o par e sai
como \ce{Br-} ($-1$); o hidrogênio se liga a um carbono, e o outro
carbono fica com uma carga positiva: \ce{CH3-CH2+} ($+1$).
\end{solution}

\begin{solution}{exo:g12:curly-arrows:8}
1.~\ce{CH3-CH2-OH + H+ -> CH3-CH2-OH2+}: doador, um \omterm{def:g10:lewis-and-shape:lone-pair}{par não ligante} do
oxigênio; aceptor, o \ce{H+}. 2.~\ce{CH3-CH2-OH2+ -> CH3-CH2+ + H2O}: o
par \ce{C-O} vai para o oxigênio positivo (aceptor), que sai com ele na
forma de água. 3.~\ce{CH3-CH2+ -> CH2=CH2 + H+}: o par de uma ligação
\ce{C-H} (doador) se desloca para o carbono positivo (aceptor) para
formar a \omterm{def:g10:lewis-and-shape:multiple-bond}{ligação dupla}, e o \ce{H+} sai. O \omterm{def:g9:ions:ion}{íon} hidrogênio usado na etapa
1 é devolvido na etapa 3: ele aparece dos dois lados e se cancela.
\end{solution}

\begin{solution}{exo:g12:curly-arrows:9}
Uma seta de um \omterm{def:g10:lewis-and-shape:lone-pair}{par não ligante} do oxigênio da água até o carbono
positivo. Depois da perda de \ce{H+}, o etanol \ce{CH3-CH2-OH}. No
total: \ce{CH2=CH2 + H2O -> CH3-CH2-OH}, uma \omterm{def:g12:curly-arrows:categories}{adição}.
\end{solution}

\begin{solution}{exo:g12:curly-arrows:10}
Dois pares se deslocam, numa só etapa: a ligação \ce{C-O} se forma, a
ligação \ce{C-Br} se rompe. Nenhum intermediário.
\end{solution}

\begin{solution}{exo:g12:curly-arrows:11}
O cloro é mais eletronegativo ($3.16 - 2.55 = 0.61$ contra
$2.96 - 2.55 = 0.41$): o carbono do clorometano é o mais $\delta^+$, o
que sozinho o faria reagir mais depressa. Não é o que acontece: a
velocidade depende também da facilidade com que o \omterm{def:g10:electron-shells:family}{halogênio} sai com o
par, e o bromo, maior, sai com mais facilidade.
\end{solution}

\begin{solution}{exo:g12:curly-arrows:12}
Se o hidrogênio se liga ao carbono 1: \ce{CH3-CH+-CH3}, que dá o
2-cloropropano \ce{CH3-CHCl-CH3}. Se ele se liga ao carbono 2:
\ce{+CH2-CH2-CH3}, que dá o 1-cloropropano \ce{CH2Cl-CH2-CH3}.
\end{solution}

\begin{solution}{exo:g12:curly-arrows:13}
A seta vai no sentido errado: ela tem de partir do doador, um \omterm{def:g10:lewis-and-shape:lone-pair}{par não
ligante} do \ce{OH-}, e terminar no carbono aceptor. E a nova ligação é
entre o oxigênio e o carbono, não entre o oxigênio e o bromo: o bromo
sai, levando o par da ligação \ce{C-Br}.
\end{solution}

\begin{solution}{exo:g12:curly-arrows:14}
O carbono do grupo \ce{C=O} de um ácido está ligado a dois \omterm{def:g7:atoms-and-molecules:atom}{átomos} de
oxigênio eletronegativos: fortemente $\delta^+$, um aceptor. O oxigênio
do \omterm{def:g11:functional-groups:alcohol}{álcool} tem $\delta^-$ e dois \omterm{def:g10:lewis-and-shape:lone-pair}{pares não ligantes}: um doador. No total,
perde-se água; o \ce{-OH} do ácido é trocado pelo \ce{-O-R} do \omterm{def:g11:functional-groups:alcohol}{álcool}:
uma \omterm{def:g12:curly-arrows:categories}{substituição}.
\end{solution}

\begin{solution}{exo:g12:curly-arrows:15}
Empurrado pelos \omterm{def:g9:inside-the-atom:nucleus}{elétrons} da \omterm{def:g10:lewis-and-shape:multiple-bond}{ligação dupla}, o par da ligação \ce{Br-Br}
se desloca para o bromo mais distante: o bromo mais próximo fica
$\delta^+$, um sítio aceptor. Primeira etapa: uma seta da \omterm{def:g10:lewis-and-shape:multiple-bond}{ligação dupla}
até o bromo mais próximo, e uma da ligação \ce{Br-Br} até o bromo mais
distante, que sai como \ce{Br-}; o outro carbono da antiga \omterm{def:g10:lewis-and-shape:multiple-bond}{ligação dupla}
fica com uma carga positiva.
\end{solution}

\begin{solution}{pb:g12:curly-arrows:1}
\textbf{1.} Eteno: \ce{H2C=CH2}, uma \omterm{def:g10:lewis-and-shape:multiple-bond}{ligação dupla} e quatro \ce{C-H}.
Brometo de hidrogênio: \ce{H-Br} com três \omterm{def:g10:lewis-and-shape:lone-pair}{pares não ligantes} no bromo.

\textbf{2.} A \omterm{def:g10:lewis-and-shape:multiple-bond}{ligação dupla}: o seu segundo par é uma região rica em
\omterm{def:g9:inside-the-atom:nucleus}{elétrons}.

\textbf{3.} $2.96 - 2.20 = 0.76$: o \ce{H} tem $\delta^+$.

\textbf{4.} O \omterm{def:g7:atoms-and-molecules:atom}{átomo} de hidrogênio.

\textbf{5.} Uma seta da \omterm{def:g10:lewis-and-shape:multiple-bond}{ligação dupla} até o hidrogênio do \ce{HBr}; uma
seta da ligação \ce{H-Br} até o bromo.

\textbf{6.} O carbocátion \ce{CH3-CH2+} ($+1$) e o \ce{Br-} ($-1$).

\textbf{7.} Antes: $0 + 0 = 0$; depois: $+1 - 1 = 0$.

\textbf{8.} Uma seta de um \omterm{def:g10:lewis-and-shape:lone-pair}{par não ligante} do \ce{Br-} até o carbono
positivo.

\textbf{9.} Seis (três ligações): falta um par para o octeto, por isso
ele aceita um par de qualquer doador próximo.

\textbf{10.} O carbocátion \ce{CH3-CH2+}.

\textbf{11.} Ele é formado na primeira etapa e consumido na segunda.

\textbf{12.} Uma \omterm{def:g12:curly-arrows:categories}{adição}.

\textbf{13.} Só o grupo: a cadeia de dois carbonos é conservada, e a
\omterm{def:g10:lewis-and-shape:multiple-bond}{ligação dupla} vira uma ligação simples que carrega \ce{H} e \ce{Br}.

\textbf{14.} \qty{28.0}{g/mol}; \qty{108.9}{g/mol}.

\textbf{15.} $10.0 / 28.0 = \qty{0.357}{mol}$.

\textbf{16.} Eteno: $0.357 - x$; \ce{HBr}: em excesso; bromoetano: $x$.
O eteno é o limitante: $x_{\max} = \qty{0.357}{mol}$.

\textbf{17.} $0.357 \times 108.9 = \qty{38.9}{g}$.

\textbf{18.} $0.75 \times 38.9 = \qty{29.2}{g}$.

\textbf{19.} Cerca de \qty{29}{g} de bromoetano.
\end{solution}
